\section{An exact-shell density theorem at fixed rank and inner partition}
\label{sec:fixed-rank-density}

The density population of SH(b) allows the rank and all three
partitions to vary.  The following restricted theorem uses its
literal sum-of-binary-lengths encoding and its nonzero initial
coefficient conditioning, but fixes the rank and inner partition.
It does not substitute a Euclidean-size population for that encoding.

\begin{theorem}[PNEW-009: fixed proper embeddings have exponentially sparse strict saturation exceptions]
Fix $n\ge2$ and a polynomial highest weight $\mu$ of $GL_n$.
Put $d=|\mu|$ and $m=\dim S_\mu\mathbb C^n$, and suppose
$m>n$.  Among all ordinary plethysm inputs $(\lambda,\mu,\pi)$
with $\operatorname{ht}(\lambda)\le m$, $\operatorname{ht}(\pi)\le n$,
$a_{\lambda,\mu}^{\pi}>0$ and
\[
              \ell(\lambda)+\ell(\mu)+\ell(\pi)=N,
\]
the proportion with $s(a_{s\lambda,\mu}^{s\pi})>0$ is
$O_{n,\mu}(N^{m+n-1}2^{-\eta_{n,\mu}N})$ for some
$\eta_{n,\mu}>0$.  The denominator is nonzero for every
sufficiently large $N$.

In particular the narrow SH(b) conclusion holds on every fixed
rank, fixed cubic-inner-partition population with $n\ge2$.
No uniform estimate on the constants as $n$ or $\mu$ varies is
asserted.
\end{theorem}

We isolate the elementary estimates and the classical branching
input, so the scope of the conclusion can be checked separately.

\begin{lemma}[Polynomial sublevel estimate]
Let $P$ be a nonzero homogeneous polynomial of degree $e\ge1$
on $\mathbb R^r$, and fix a full lattice $\Lambda$ there.
For each fixed $A,C>0$, the number of $u\in\Lambda$ with
\[
       \|u\|\le CR,\qquad |P(u)|\le AR^{e-1}
\]
is $O(R^{r-1/e})$ as $R\to\infty$.
\end{lemma}
\begin{proof}
Choose a primitive lattice vector $v$ with $P(v)\ne0$;
such a vector exists because a nonzero polynomial cannot vanish
on every lattice point.  Complete $v$ to a lattice basis and
write its coordinates as $(y_1,\ldots,y_r)$.  For fixed
$y_2,\ldots,y_r$, the polynomial in $y_1$ has degree $e$
and constant leading coefficient $P(v)$.  Factor it over
$\mathbb C$.  If its absolute value is at most $AR^{e-1}$,
at least one of its $e$ linear factors has absolute value at
most $(A/|P(v)|)^{1/e}R^{1-1/e}$.  Its real arguments then
lie in at most $e$ intervals of that radius, irrespective
of the imaginary parts of the roots.  There are
$O(R^{1-1/e})$ integral choices for $y_1$ and $O(R^{r-1})$
choices for the remaining coordinates.  A fixed lattice basis
change changes only the constants in the norm bound.
\end{proof}

\begin{lemma}[Piecewise quasipolynomials and strict saturation]
Let $C\subset E\simeq\mathbb R^r$ be a full-dimensional
rational polyhedral cone, with lattice $\Lambda$.
Suppose $M:C\cap\Lambda\to\mathbb Z_{\ge0}$ is given by
finitely many quasipolynomials on a finite conic polyhedral
subdivision.  Assume that on every ray the resulting stretching
function is interpreted by these global constituent polynomials.
Let $e_*$ be the maximum of 1 and the degrees of the nonzero
constituent polynomials on full-dimensional chambers.  Then
\[
 \#\{u\in C\cap\Lambda:\|u\|\le R,
                      \ s(M(su))>0\}
                 =O(R^{r-1/e_*}).
 \tag{PNEW-density-Euclidean}
\]
The same estimate holds when one conditions on $M(u)>0$.
\end{lemma}
\begin{proof}
Discard the finitely many chamber walls at a cost of
$O(R^{r-1})$ lattice points.  Fix a full-dimensional chamber
and a common finite-index period lattice for all its
quasipolynomials.  Write its constituent polynomials as
$P_\rho=\sum_{j=0}^{e_\rho}P_{\rho,j}$ with homogeneous
pieces.  If $P_\rho\not\equiv0$, its top piece is nonnegative
on the chamber.  Indeed, for any point $v$ in its interior,
choose lattice points in the residue class $\rho$ within
bounded distance of $tv$.  They remain in the chamber for
large $t$, and
$t^{-e_\rho}P_\rho(tv+O(1))\to P_{\rho,e_\rho}(v)$;
the terms before the limit are nonnegative.  Continuity gives
the assertion on the closure as well.

For $\|u\|\le R$ and $e_\rho\ge1$, the sum of the absolute
values of all lower pieces at $u$ is at most
$A_\rho R^{e_\rho-1}$.  If
\[
        P_{\rho,e_\rho}(u)>A_\rho R^{e_\rho-1},
 \tag{PNEW-leading-dominates}
\]
then $P_\rho(Xu)>0$ for every real $X\ge1$: multiply the
top term by $X^{e_\rho}$, and bound every lower power by
that same power.  A nonzero degree-zero constituent is a
positive constant; a zero polynomial is ignored by the
definition of the strict saturation index.

The constituents of $s\mapsto M(su)$ are among the
polynomials $X\mapsto P_\rho(Xu)$: choose a common scalar
period of the finite quotient lattice and restrict to each
residue of $s$.  Thus all nonzero constituents are positive
at every positive integer whenever
(PNEW-leading-dominates) holds for every applicable nonzero
$P_\rho$.  Each failure lies in a polynomial sublevel set
covered by the preceding lemma.  There are only finitely
many chambers and residue polynomials.  Summing their
bounds and the wall contribution proves the estimate.
Conditioning on $M(u)>0$ can only decrease this numerator.
\end{proof}

\begin{lemma}[From size bounds to literal binary-length shells]
Let $D\ge3$ and
\[
 E=\{(x_1,\ldots,x_D):c_1x_1+\cdots+c_Dx_D=0\},
 \qquad c_i\in\mathbb Z\setminus\{0\},
\]
where coefficients of both signs occur.  Let $\Lambda=E\cap
\mathbb Z^D$, so $r=\dim E=D-1$.  Suppose a nonnegative
piecewise quasipolynomial $M$ on a finite union of rational
cones in $E\cap\mathbb R_{\ge0}^D$ has a full-dimensional
chamber with a nonzero residue polynomial and with strictly
positive coordinates.  Suppose its exceptional lattice
points of norm at most $R$ number $O(R^{r-\epsilon})$,
where $0<\epsilon\le1$.

Allow zero coordinates and omit them when summing binary
lengths.  Then among the points with $M(u)>0$ and
$\sum_i\ell(u_i)=L$, the exceptional proportion is
\[
        O\!\left(L^{D-1}2^{-\epsilon L/(2D(D-1))}\right),
 \tag{PNEW-raw-shell-bound}
\]
provided exceptions on proper coordinate faces are counted
arbitrarily.  The denominator is nonzero for every sufficiently
large $L$.
\end{lemma}
\begin{proof}
\emph{A denominator in every large shell.}
Choose a positive vector $v$ in the interior of a chamber
where a constituent has strictly positive top homogeneous
piece.  Avoid all equalities $v_i=2^k v_j$.
Also avoid dyadic values of its individual coordinates, by
a sufficiently small scalar adjustment if necessary.
Only finitely many of these hyperplanes meet a fixed small
compact positive neighborhood.  None contains $E$ when
$D\ge3$ and every $c_i$ is nonzero, so the choice is possible.
As $t$ increases from 1 to 2, each coordinate of $tv$ crosses
one binary threshold, and these $D$ crossings are distinct.
There are consequently positive open subintervals on which
the sum of its coordinate binary lengths takes $D$
consecutive integer values $C_0,\ldots,C_0+D-1$.
Choose a small open ball $B_j\subset E$ around one point
of each such interval, with the same length pattern and
positive top piece throughout its closure.

On $2^bB_j$ the total length is $Db+C_0+j$.
In a fixed nonzero residue class the constituent is positive
there for all large $b$, uniformly on the ball; its top
piece dominates its lower pieces.  The number of points of
that full-rank lattice coset in $2^bB_j$ is
$\Omega(2^{br})$ (inscribe a fixed parallelepiped and count
its scaled lattice translates).  The $D$ consecutive values
of $C_0+j$ cover every residue modulo $D$.  Therefore the
number of nonzero inputs in every sufficiently large shell
is at least
\[
                         c\,2^{rL/D}
 \tag{PNEW-shell-denominator}
\]
for one fixed $c>0$.

\emph{Unbalanced sizes and omitted coordinates.}
First suppose all $D$ coordinates are positive.  Fix their
lengths $b_1,\ldots,b_D$ with sum $L$, and let
$B=\max b_i$.  The equation for $E$ determines a coordinate
having length $B$ uniquely from all the others.  Thus the
number of choices is at most $2^{L-B}$, whether or not the
determined value is an integer.  There are at most
$D(L+1)^{D-1}$ possible length patterns and choices of a
maximum coordinate.  Put
\[
                      \delta=\frac{\epsilon}{2D(D-1)}.
\]
Those with $B>L/D+\delta L$ therefore number
$O(L^{D-1}2^{rL/D-\delta L})$.

If precisely $D'<D$ coordinates are positive, the same
argument gives $O(L^{D'-1}2^{(1-1/D')L})$, unless all
remaining coefficients have one sign, in which case there
are no such points.  The possible empty or single-coordinate
cases contribute no large shell.  Since
\[
 \frac{D-1}{D}-\frac{D'-1}{D'}\ge\frac1{D(D-1)}\ge\delta,
\]
all proper coordinate faces have the required exponentially
smaller count, even if every such point is exceptional.

\emph{Balanced exceptional points.}
For the remaining full-height inputs,
$\|u\|\le C_D2^{L/D+\delta L}$, so the assumed size bound
gives at most
\[
 O\!\left(2^{(r-\epsilon)(L/D+\delta L)}\right)
       \le O\!\left(2^{rL/D-\epsilon L/(2D)}\right).
\]
Combine the three numerator bounds and divide by
(PNEW-shell-denominator).  This proves
(PNEW-raw-shell-bound) without replacing its length shell
by a norm ball.
\end{proof}

\paragraph{The classical branching input and the full cone.}
The restriction associated to the fixed inner representation
is a fixed homomorphism $SL_n\to SL_m$ with finite kernel.
For the compact branching theorem take $K$ to be the image
of $SU(n)$ in $SU(m)$.  Its weight lattice identifies with
a finite-index sublattice of the $SU(n)$ weight lattice.
Apply the theorem on that image lattice, and extend the
multiplicity function by zero on the other congruence
classes.  This remains a finite piecewise quasipolynomial
description on the full lattice.
The map of Lie algebras is injective: $\mathfrak{sl}_n$ is
simple and the representation is nontrivial.  Its image is
proper since $m>n$.  No nonzero ideal of this image is an
ideal of $\mathfrak{sl}_m$, because the latter is simple.
The standard full-cone lemma for branching therefore gives
a solid cone in the product of the two semisimple weight
spaces, of dimension $(m-1)+(n-1)$.
A precise accessible statement is Baldoni and Vergne (2015),
Section 5.1, Lemma 61.
Their Theorem 63 and Propositions 65--67 supply a finite
conic piecewise quasipolynomial formula for branching
multiplicities, including global stretching constituents.
These are classical branching theorems used as cited inputs;
we do not claim to reprove the quantization theorem from
which the chamber description follows.

For clarity, passing to polynomial $GL$ weights introduces
exactly the degree equation
\[
                      d|\lambda|=|\pi|.
 \tag{PNEW-degree-hyperplane}
\]
Every semisimple branching pair lifts to this hyperplane on
its appropriate finite-index congruence lattice.  Indeed the
central $n$-th roots acting on $S_\mu\mathbb C^n$ have
character $z^d$, which gives the necessary congruence for
the determinant exponent of the target.  Adding one to all
$m$ parts of $\lambda$ adds $e=dm/n$ to all $n$ parts of
$\pi$.  The number $e$ is an integer because
$\det(S_\mu\mathbb C^n)=(\det\mathbb C^n)^{dm/n}$:
the total exponent $dm$ is equally distributed among the
$n$ variables.  This determinant direction is independent
of the semisimple weight coordinates.  Adding a sufficiently
large multiple makes every part positive.  Thus the
polynomial branching cone has dimension $m+n-1$, is full
inside (PNEW-degree-hyperplane), and has an open set of
positive, strictly dominant weights.

\begin{proof}[Proof of PNEW-009]
Apply the piecewise quasipolynomial lemma to the branching
function in the degree hyperplane, with $D=m+n$,
$r=D-1$ and $\epsilon=1/e_*$.  A full-dimensional chamber
has some nonzero residue polynomial because its cone is
the asymptotic support of branching.  The preceding
paragraph ensures a chamber of positive polynomial weights.
Partition dominance imposes linear inequalities only and
does not affect either lattice estimate.

Apply the binary-length lemma with $m$ coefficients equal
to $d$ and $n$ coefficients equal to $-1$.
It uses $L=N-\ell(\mu)$, where $\ell(\mu)$ is a fixed
constant.  Thus one may take
\[
                \eta_{n,\mu}
                    =\frac1{2(m+n)(m+n-1)e_*}
\]
in the displayed polynomial-times-exponential bound,
absorbing the fixed length offset in the implicit constant.
The numerator is precisely the positive-initial inputs
with a nonzero strict saturation index; the denominator is
the positive-initial population of the stated shell.
For $\mu=(3)$ one has $m=\binom{n+2}{3}>n$ for every
$n\ge2$, so the final assertion follows.
\end{proof}

\paragraph{What remains for the original SH(b).}
PNEW-009 eliminates a fixed-embedding density obstruction,
but its degree bound, periods, chamber constants and onset
are allowed to depend on $n$ and $\mu$.  The original
population sums over changing embeddings.  A countable
union of exponentially sparse fixed-parameter families need
not be sparse in that growing population.  The remaining
gate is a uniform estimate strong enough to sum their
exceptional counts, or an explicit family violating such
an estimate.  No such uniform estimate is proved here.

